% !TeX program = xelatex
\documentclass[11pt]{article}

\usepackage[a4paper, margin=1in]{geometry}
\usepackage{amsmath, amssymb, amsthm}
\usepackage{stmaryrd}
\usepackage{listings}
\usepackage{xcolor}
\usepackage{hyperref}
\usepackage{microtype}
\usepackage{fontspec}
\usepackage{booktabs}
\setmonofont{DejaVu Sans Mono}[Scale=0.85]

\sloppy
\setlength{\emergencystretch}{3em}

\newcommand{\btt}[1]{\texttt{\detokenize{#1}}}
\newcommand{\bv}[1]{\llbracket #1 \rrbracket}
\newcommand{\feq}{\mathrel{=_{[F]}}}

\theoremstyle{definition}
\newtheorem*{principle}{Principle}
\newtheorem*{observation}{Observation}

% ---- Lean code listing style -------------------------------------------
\lstdefinelanguage{lean}{
  morekeywords={theorem, lemma, def, by, fun, let, in, do, match, with,
    if, then, else, return, import, open, namespace, end, structure,
    inductive, partial, private, noncomputable, syntax, elab, elab_rules,
    macro, macro_rules, tactic, conv, command, example, have, show, intro,
    apply, simp, rw, ring, all_goals, repeat, decide, native_decide,
    fin_cases, norm_num, exact, set_option, refine},
  sensitive=true,
  morecomment=[l]{--},
  morecomment=[s]{/-}{-/},
  morestring=[b]"
}
\lstset{
  inputencoding=utf8,
  language=lean,
  basicstyle=\ttfamily\footnotesize,
  keywordstyle=\color{blue!70!black}\bfseries,
  commentstyle=\color{gray}\itshape,
  stringstyle=\color{red!60!black},
  showstringspaces=false,
  breaklines=true,
  columns=fullflexible,
  frame=single,
  framesep=4pt,
  rulecolor=\color{black!30},
  backgroundcolor=\color{gray!6},
  xleftmargin=6pt,
  xrightmargin=6pt,
  aboveskip=0.8\baselineskip,
  belowskip=0.8\baselineskip
}
% -----------------------------------------------------------------------

\title{The Forbid-Free Flag Framework:\\ Design and Rationale}
\author{}
\date{}

\begin{document}
\maketitle

\begin{abstract}
This note explains the \emph{forbid-free} layer of the Lean flag-algebra
development: the machinery that, when a proof works under a forbidden-subgraph
hypothesis (triangle-free, $K_4$-free, $\dots$), generates and reasons with only
the flags that survive that hypothesis, rather than the full flag enumeration.
It is meant to be read on its own. It first recalls the small amount of
flag-algebra background it needs, then develops the design in the order the
pieces depend on one another, and along the way records the observations and
design decisions that shaped it---in particular \emph{where} the speed-up
actually comes from, why the flag indexing is what it is, and the trade-offs of
a more aggressive ``true pruning'' that we deliberately did not build.
Identifier names follow the repository convention \btt{_n_k_m_i}: $n$ vertices,
type $\sigma$ the $k$-vertex graph of index $m$ (\btt{_0_0} is the empty type
$\emptyset_t$), enumeration index $i$.
\end{abstract}

\tableofcontents

\section{Background and motivation}\label{sec:motivation}

A flag-algebra density proof works by expanding the quantity of interest in a
basis of \emph{flags} of some fixed size $n$, and then bounding the resulting
linear (or, with a semidefinite certificate, quadratic) combination. The
repository realises this concretely: for each flag it generates a named Lean
constant \btt{Flag_n_k_m_i} (and its quotient \btt{FlagAlgebra_n_k_m_i}),
together with the bridging lemmas a proof needs---how the flag unlabels, how it
projects downward, how products of smaller flags expand into it. All of this is
produced by elaboration-time \emph{generation commands} (\btt{generate_flags},
\btt{generate_empty_typed_flags}) that enumerate the isomorphism classes,
deduplicate them, and emit the constants with machine-checked bridging lemmas.

The number of flags grows quickly: there are $11$ unlabeled graphs on $4$
vertices, $34$ on $5$, $156$ on $6$, $1044$ on $7$, and the labeled (typed)
counts are larger still. Generating every flag---each with a kernel-checked
\btt{type_embed} proof, a \btt{downward} lemma, and an entry in the completeness
argument---is the dominant cost of setting up a proof, and at $n \ge 6$ it is
the difference between a proof that compiles and one that does not.

Many extremal problems are stated under a \emph{forbidden subgraph}: maximise
edge density among triangle-free graphs (Mantel), among $K_r$-free graphs
(Tur\'an), and so on. Under such a hypothesis the flags that \emph{contain} the
forbidden graph $H$ carry no information: they are identically zero in the
relevant conditioned algebra. So in principle we never need them. The
forbid-free framework makes that principle operational:

\begin{principle}
Under a forbidden subgraph $H$, generate only the $H$-free flags, prove they are
\emph{complete} (they account for everything that survives the forbidden
relation), and route every proof bridge through that smaller set.
\end{principle}

The rest of this note is about turning that one sentence into something that
both type-checks and is actually faster.

\section{The vanishing principle}\label{sec:vanishing}

The semantic foundation already existed in \btt{Forbid/Basic.lean} before the
forbid-free layer was built; we recall it because everything downstream keys off
one predicate it produces.

Write $f \feq g$ (\btt{forbidEq}, notation \btt{=[F]}) for the relation ``$f$ and
$g$ agree almost surely under random homomorphisms drawn conditioned on the
forbidden flag $F$ having density zero.'' It is an additive, scalar-linear
congruence. The key fact is that a flag which \emph{does} contain $F$ dies:

\begin{lstlisting}
theorem basisVector_forbidEq_zero (F_forbid : FinFlag emptyT) (F : FinFlag sigma)
    (hF : flagDensity1 F_forbid.2 (unlabel F.2) > 0) : llbracket basisVector F rrbracket =[F_forbid] 0
\end{lstlisting}

Feeding this into the ordinary size-$\ell$ expansion of a flag (resp.\ a product
of flags) and discarding the vanishing terms yields the two lemmas the whole
framework is built to consume:
\[
  \bv{\mathsf{bv}\,F}
  \feq
  \!\!\sum_{\substack{F' : \mathrm{FlagWithSize}\,\sigma\,\ell \\[1pt]
        \mathrm{flagDensity}_1\, F\, (\mathrm{unlabel}\,F') = 0}}\!\!
        \bigl(\mathrm{flagDensity}_1\, F.2\, F'\bigr)\cdot \bv{\mathsf{bv}\,\langle\ell,F'\rangle},
\]
and analogously \btt{basisVector_quot_mul_forbidEq_sum} for a product, with
$\mathrm{flagDensity}_2$ in place of $\mathrm{flagDensity}_1$. The sum ranges
\emph{only} over the flags $F'$ whose forbidden-density is zero.

\begin{observation}[The linchpin predicate]
The index set of these expansions is
\[
   \mathrm{univ}.\mathrm{filter}\;\bigl(\lambda F'.\;
     \mathrm{flagDensity}_1\, F_{\mathrm{forbid}}.2\,(\mathrm{unlabel}\,F') = 0\bigr).
\]
Call the body the \emph{forbid-free predicate} $p$. It is a density-based,
\emph{analytic} statement (``$F$ occurs in $F'$ with density $0$''), not an
\emph{a priori} combinatorial one. Every design decision below is organised so
that the generated $H$-free flag set is provably equal to exactly this
$\mathrm{univ}.\mathrm{filter}\,p$, because that is what lets a bridge rewrite an
expansion's sum directly onto the named generated constants.
\end{observation}

\section{Forbid-free generation and filtered completeness}\label{sec:generation}

\paragraph{The first design decision: stay in the full algebra.}
We do \emph{not} build a separate ``$H$-free flag algebra'' with its own type,
operations, and metatheory. We stay inside the existing \btt{FlagAlgebra sigma}
and the relation $\feq$, and change only \emph{which flag constants we bother to
generate}. The cost of a separate algebra (re-proving the entire flag-algebra
toolkit over it) would dwarf any generation savings, and---more importantly---it
would not compose with the existing bridges and certificates. So the forbid-free
layer is additive: it sits on top of the ordinary flags.

\paragraph{The completeness we must prove.}
For the generated $H$-free set \btt{flagSetHfree} to be usable by the bridges, it
must be proved equal to the linchpin set, in exactly that form:
\begin{lstlisting}
theorem flagSetHfree_n_k_m_F_eq :
    flagSetHfree_n_k_m_F = Finset.univ.filter (fun F' => flagDensity1 F.toFinFlag.2 (unlabel F') = 0)
\end{lstlisting}
Matching the predicate verbatim is the whole point: a bridge can then take the
expansion of Section~\ref{sec:vanishing}, whose sum is over
$\mathrm{univ}.\mathrm{filter}\,p$, and \btt{rw} it onto \btt{flagSetHfree}.

\paragraph{The commands.}
Two commands generate the constants and prove this completeness:
\begin{itemize}
  \item \btt{generate_forbid_free_empty_typed_flags n Forbid} --- the empty-typed
        ($k=m=0$) flags on $n$ vertices;
  \item \btt{generate_forbid_free_flags n k m Forbid} --- the $\sigma$-typed
        analogue (signature leads with the flag size $n$, matching
        \btt{generate_flags n k m}).
\end{itemize}
Each emits only the $H$-free flag constants, their \btt{unlabel}/\btt{downward}
bridges, the set \btt{flagSetHfree}, the completeness lemma above, and a
\btt{flagSetHfree_..._val_eq} stating that the set's underlying multiset is the
explicit list of named free constants (the bridges need both: \btt{_eq} to
recognise the set, \btt{_val_eq} to expand a sum over it into named terms).

\paragraph{How completeness is proved (Phase 1).}
The generator already proves the \emph{full} completeness
$\btt{flagSet} = \mathrm{univ}$ (\btt{genFlagSet_eq_univ}). The forbid-free
completeness is then obtained by \emph{filtering} that:
\begin{enumerate}
  \item a single \btt{native_decide} checks that the full Lean enumeration,
        filtered by an executable $H$-freeness test, equals the explicit list of
        free constants;
  \item filtering commutes with \btt{toFinset} and with \btt{= univ}, giving
        $\btt{flagSetHfree} = \mathrm{univ}.\mathrm{filter}\,p$.
\end{enumerate}
The executable $H$-freeness test routes through the computable
\btt{sym2EmptyTypeFlagDensity1} (the noncomputable \btt{flagDensity1} cannot be
run by \btt{native_decide} directly) and is reconciled with the analytic
predicate by \btt{flagDensity1_eq_sym2EmptyTypeFlagDensity1}.

\paragraph{Typed vs.\ empty-typed, and the \texttt{toUnderlying} bridge.}
$H$-freeness is a property of the \emph{underlying graph}: a $\sigma$-typed flag
is $H$-free iff the graph obtained by forgetting its labelling is. To express
this we added a computable forgetful map and its compatibility lemma:
\begin{lstlisting}
def Sym2Flag.toUnderlying (S : Sym2Flag sigma n) : Sym2EmptyTypedFlag n
theorem Sym2Flag.unlabel_toFlag_eq (S : Sym2Flag sigma n) :
    unlabel (S.toFlag) = (S.toUnderlying).toFlag
\end{lstlisting}
With these, the typed forbid-free predicate is decided on the underlying graph,
and the typed completeness reuses the empty-typed reasoning. (At a size where the
ordinary \btt{FlagDef} already generated every flag, \btt{elabUnlessDefined}
makes the forbid-free command \emph{reuse} those constants rather than redefine
them; see Section~\ref{sec:indexing}.)

\section{The forbid bridges}\label{sec:bridges}

A flag-algebra proof expands flag-algebra elements in three places, and each had
a bridge that built the \emph{full} flag set and dropped the forbidden terms with
\btt{simp}. We reworked all three to consume \btt{flagSetHfree} directly:

\begin{center}
\begin{tabular}{lll}
\toprule
What expands & Ordinary bridge & Forbid-free bridge \\
\midrule
product of flags & \btt{generate_mul_theorems} & \btt{generate_forbid_free_mul_theorems} (command)\\
the unit $1$     & \btt{expand_one_at}         & \btt{expand_one_hfree_at} (tactic)\\
a single flag    & \btt{flag_expand}           & \btt{flag_expand_hfree} (tactic)\\
\bottomrule
\end{tabular}
\end{center}

Each forbid-free bridge follows the same skeleton: apply the relevant restricted
expansion from Section~\ref{sec:vanishing}, rewrite its sum onto
\btt{flagSetHfree}, expand to the named constants via \btt{_val_eq}, and close.

\paragraph{A subtlety worth recording: decidability instances.}
The naive step ``\btt{rw [<- flagSetHfree_eq]} to turn the expansion's
$\mathrm{univ}.\mathrm{filter}\,p$ into \btt{flagSetHfree}'' \emph{fails}, and the
reason is instructive. The filter in the expansion lemma and the filter in
\btt{flagSetHfree_eq} carry different \btt{DecidablePred} instances (one arises
from an abstract $F_{\mathrm{forbid}}$ later specialised, the other is synthesised
for the concrete forbidden graph). \btt{Decidable} is \emph{data}, not a
proof-irrelevant \btt{Prop}; two instances are propositionally equal
(\btt{Subsingleton}) but not definitionally so, so \btt{rw}---which matches up to
syntactic/defeq equality---does not see them as the same term. The robust pattern is
to rewrite the \emph{set} with a congruence that discharges the instance gap:
\begin{lstlisting}
rw [Finset.sum_congr (s2 := flagSetHfree) (by rw [flagSetHfree_eq]; try congr 1) (fun _ _ => rfl)]
simp only [Finset.sum_eq_multiset_sum, flagSetHfree_val_eq]
simp                       -- evaluates the list and the (simp-tagged) density lemmas
exact forbidEq_refl _ _
\end{lstlisting}
Here \btt{congr 1} reduces the residual to the instance equality, which closes by
subsingleton; \btt{try} is needed because in the empty-typed case the instances
\emph{do} coincide and the \btt{rw} already closes the goal, so \btt{congr 1}
would error on no goals. (Two further bridge-implementation gotchas, both fixed:
the generation commands emit unqualified names that pick up the surrounding
\btt{namespace}, so cross-command ``does this lemma exist?'' checks must look up
both the current namespace and the root; and \btt{Finset.sum_eq_multiset_sum}
must be applied with \btt{simp only}, not \btt{rw}, because its higher-order
pattern $\sum_{x}f\,x$ defeats \btt{rw}'s matcher.)

\section{The completeness wall, and pruned generation}\label{sec:pruned}

\paragraph{The wall.}
Phase 1's completeness \btt{native_decide} reduces the \emph{full} enumeration.
For empty-typed flags this is fine up to $n=6$ (the enumeration has $156$
elements; it runs in under a minute). For \emph{typed} flags at $n=6$ it is not:
the typed enumeration \btt{genFlagsOrdered sigma 6} has on the order of $900$
labeled graphs before deduplication, and the \btt{native_decide} goal
\btt{(genFlagsOrdered sigma 6).filter isHfree = [free list]} forces the elaborator
to manipulate that whole object as a term. It times out at the \btt{isDefEq}
heartbeat limit after roughly $41$ minutes. Phase 1 filtering cannot avoid this:
proving ``the free list equals the full enumeration filtered'' inherently names
the full enumeration.

\paragraph{The insight that breaks the wall.}
Two facts combine. First, \emph{a typed flag is $H$-free iff its underlying graph
is}. Second, the generator already builds typed flags as
\btt{labeledOfGraph} applied to the \emph{underlying-graph} enumeration
\btt{genSym2Graphs n}. So instead of building all typed flags and filtering, we
can filter at the cheap \emph{graph} level---\btt{genSym2Graphs 6} has only $156$
elements, of which $38$ are $K_3$-free---and build typed flags only over the
surviving $38$ graphs. The expensive labeled deduplication then runs over $\sim
111$ flags instead of $\sim 900$, and the corresponding \btt{native_decide}
finishes in about two minutes.

\paragraph{The construction and its completeness.}
We added (in \btt{ForbidFreeGenerator.lean}) a generic pruned generator and its
completeness theorem:
\begin{lstlisting}
def genLabeledGraphsHfree (sigma) (n) (q : Sym2Graph n -> Bool) : List (Sym2LabeledGraph sigma n)
  -- (genSym2Graphs n).filter q  |> flatMap (labeledOfGraph sigma) |> keyed-dedup
def genFlagsHfree (sigma) (n) (q) : List (Sym2Flag sigma n)

theorem genFlagsHfree_toFinset_eq (q : Sym2Graph n -> Bool) (p : Sym2Flag sigma n -> Bool)
    (hq : forall {G G'}, G ~sf G' -> q G = q G')                      -- q iso-invariant
    (hcompat : forall Glab, p (mk Glab) = q (underlying Glab)) :       -- p agrees with q on underlyings
    (genFlagsHfree sigma n q).toFinset = Finset.univ.filter (fun F => p F = true)
\end{lstlisting}
The completeness is a membership argument that mirrors the existing
\btt{genFlagSet_eq_univ}, threading the predicate through the \emph{existing}
\btt{genSym2Graphs_complete} (every graph is isomorphic to a generated
representative). It reuses the existing deduplication-completeness lemmas and
needs only two new small facts (\btt{foldl_dedupStepL_subset}: dedup survivors
come from the input; \btt{mem_labeledOfGraph_underlying}: a labeled graph built
over $G$ has underlying graph $G$).

\begin{observation}[What we deliberately avoided]
A ``textbook'' pruned generation would never enumerate the forbidden graphs at
all, by augmenting only $H$-free representatives. Its completeness would require
a \emph{vertex-deletion monotonicity} lemma (deleting a vertex from an $H$-free
graph keeps it $H$-free), an augmentation recursion restricted to $H$-free
classes, and a bridge between combinatorial non-containment and the analytic
density-zero predicate. The graph-level-filter design needs \emph{none} of these:
$q$ is the density predicate itself, its iso-invariance is structural (it factors
through the quotient $\bv{\cdot}$), and completeness reuses the unrestricted
\btt{genSym2Graphs_complete}. This is why the framework filters rather than truly
prunes; Section~\ref{sec:truepruning} discusses the trade-off.
\end{observation}

The typed command's completeness was then re-wired to use this: a tractable
\btt{native_decide} establishes
\btt{sym2FlagSetHfree = (genFlagsHfree sigma n q).toFinset}, and
\btt{genFlagsHfree_toFinset_eq} supplies the rest (with \btt{hq} discharged by
\btt{Quotient.sound} and \btt{hcompat} by \btt{rfl} through \btt{toUnderlying}).
The full-enumeration \btt{native_decide} is gone. The empty-typed command is left
as-is, since its enumeration is small enough.

\section{Where the speed-up actually comes from}\label{sec:speedup}

This point caused genuine confusion and is worth stating carefully, because the
pruned command \emph{still enumerates every flag}.

To choose which indices are $H$-free, the command runs \btt{evalFlagDataRows k m
n}, which computes the data for \emph{all} $(k,m,n)$ flags. So the all-flags
enumeration does happen. But it happens as \emph{compiled native code producing a
plain list} that the macro reads---there is no proof obligation attached, and at
$n=6$ it costs seconds. That step was never the bottleneck.

The bottleneck was the \emph{kernel-checked completeness proof}. Doing the
enumeration as a proposition about the full \btt{genFlagsOrdered sigma 6} term
(so the elaborator and kernel must grind \btt{isDefEq}/\btt{reduceBool} over it)
is enormously more expensive than running the same algorithm as data. Pruning
moves the proof's heavy lifting onto the small filtered construction
\btt{genFlagsHfree}, and supplies the link to ``all flags'' through the
once-proved \emph{general} theorem \btt{genFlagsHfree_toFinset_eq}, so the kernel
never has to process the full typed enumeration. A secondary saving is real but
smaller: the emitted, kernel-checked \emph{declarations}---each flag's constants,
its \btt{type_embed} \btt{decide}, its \btt{downward} lemma---are produced only
for the $H$-free indices.

\begin{principle}[Slogan]
The win is not ``enumerate fewer flags as data.'' It is ``do not make the kernel
verify a proof about the full enumeration, and only emit kernel-checked
declarations for the flags you keep.''
\end{principle}

\paragraph{Honest boundary.} The one cost the current approach does \emph{not}
remove is the compiled, elaboration-time enumeration over all flags. At $n=6$
that is negligible. As $n$ grows the $H$-free fraction shrinks (roughly
$38/156$ at $n=6$, $\sim107/1044$ at $n=7$) and that full enumeration---together
with its $O(g^2)$ deduplication---could itself become the limiting factor. Only a
true pruned augmentation would remove it.

\section{The indexing scheme}\label{sec:indexing}

The index $i$ in \btt{Flag_n_k_m_i} is the flag's \emph{rank in the canonical
sort of all $(k,m,n)$ flags} (sorted by \btt{labeledFlagKey}). The forbid-free
command keeps these original indices and simply omits the forbidden ones, so the
generated free flags carry \emph{non-contiguous} indices with gaps. For example
the $K_3$-free $4$-vertex flags of type \btt{2_0} are
\[
  \btt{Flag_4_2_0_0},\,\dots,\,\btt{Flag_4_2_0_10},\,
  \btt{Flag_4_2_0_13},\,\btt{Flag_4_2_0_14},\,\btt{Flag_4_2_0_15},\,\btt{Flag_4_2_0_18}
\]
--- the indices $11,12,16,17$ are absent (those flags contain a triangle). They
are \emph{not} renumbered to $0,\dots,14$.

This is a deliberate, load-bearing choice: \btt{Flag_n_k_m_i} denotes the same
flag no matter which command emitted it. That single fact is what lets the
forbid-free layer be additive rather than a fork:
\begin{itemize}
  \item at sizes the ordinary \btt{FlagDef} already covers, the forbid-free
        command's \btt{elabUnlessDefined} finds the constants and \emph{reuses}
        them, instead of creating a clashing second definition;
  \item the per-flag \btt{unlabel} lemma refers to its underlying flag by that
        flag's (full-enumeration) index, and the density/multiplication theorems
        refer to flags by index---all of which line up only because the indexing
        is shared;
  \item the indices agree with the Flagmatic certificate numbering the rest of
        the pipeline assumes.
\end{itemize}

\section{Design discussion: ``true'' pruned augmentation}\label{sec:truepruning}

It is natural to ask whether we should go further and never generate the
forbidden flags at all. The honest analysis follows; the conclusion is that the
current ``filter, don't prune'' design is the better-engineered point on the
curve, and that true pruning is a \emph{conditional future} step.

\paragraph{It forfeits the shared indexing.}
The index $i$ is a flag's rank among \emph{all} flags. Computing that rank
requires knowing how many flags---forbidden or not---sort before it, which means
knowing the forbidden flags. A true pruned augmentation, by construction, never
produces (or even keys) the forbidden flags, so it can recover the free flags'
\emph{relative} order but not their absolute rank. Its natural output is a
\emph{contiguous, free-only} indexing, which denotes different flags than the
current scheme; and the full rank is genuinely unrecoverable without
re-enumerating the forbidden flags. So true pruning and the shared indexing of
Section~\ref{sec:indexing} are mutually exclusive.

\paragraph{Advantages of true pruning.}
\begin{itemize}
  \item It removes the one residual cost of Section~\ref{sec:speedup}: generation
        would scale with the number of $H$-free flags, not the total. This is what
        would matter at $n \ge 7$, where the free fraction is small and even the
        compiled full enumeration is heavy.
  \item It is conceptually cleaner---``generate exactly the admissible flags.''
  \item Its free-only numbering may align better with Flagmatic's own
        forbid-aware flag numbering (Flagmatic typically numbers only admissible
        flags), which could simplify certificate ingestion.
\end{itemize}

\paragraph{Disadvantages of true pruning.}
\begin{itemize}
  \item \emph{Loss of shared indexing} (above): name clashes with \btt{FlagDef},
        broken cross-references (\btt{unlabel} targets, density/mul theorems), and
        a certificate re-indexing problem. One would need either a separate
        \btt{FlagHfree_*} namespace or a consistent project-wide re-index.
  \item \emph{A much harder completeness proof}: the vertex-deletion monotonicity
        lemma, the $H$-free-restricted augmentation recursion, and the
        combinatorial-to-analytic ($H$-free $\Leftrightarrow$ density $0$) bridge
        --- exactly the machinery the current design avoids by reusing
        \btt{genSym2Graphs_complete}.
  \item \emph{More code and more proof risk} for a payoff that is not yet on the
        table.
\end{itemize}

\paragraph{Verdict.}
The indexing incompatibility is decisive rather than incidental: the current
scheme's value is precisely that the forbid-free machinery sits \emph{on top of}
the existing flags and certificates instead of forking them. Treat true pruning
as worthwhile only once (a)~$n \ge 7$ is a real target and (b)~the
elaboration-time enumeration---not the proofs---is measured to be the
bottleneck. If that day comes, the cleanest design is to let the pruned generator
renumber to free-rank under its own names and re-express the certificate in
free-rank (likely Flagmatic's native numbering anyway), rather than trying to
preserve full-enumeration indices.

\section{Status, validation, and where things live}\label{sec:status}

\paragraph{Validated.} Empty-typed generation works through $n=6$; typed
generation works through $n=6$ (the $K_3$-free type-\btt{1_0} case builds in
about $1$m$48$s, versus the $41$-minute timeout of the unpruned approach). All
three bridges are exercised end-to-end by \btt{Flagmatic/MantelHfree.lean}, a
forbid-free formalisation of Mantel's theorem that proves the $1/2$ bound using
forbid-free generation, the forbid-free multiplication theorems, an
\btt{flag_expand_hfree} objective expansion, and \btt{expand_one_hfree_at}; it
builds green and leaves the ordinary \btt{Mantel.lean} untouched. (At $n=3$ there
is no generation saving---\btt{FlagDef} already has every flag---so \btt{MantelHfree}
validates the bridges rather than the speed-up; the speed-up is what makes the
$n=6$ typed generation feasible at all.)

\paragraph{API summary.} The generation API consistently leads with the flag
size $n$:
\begin{center}
\begin{tabular}{ll}
\toprule
Command / tactic & Role \\
\midrule
\btt{generate_forbid_free_empty_typed_flags n Forbid} & empty-typed $H$-free flags + completeness \\
\btt{generate_forbid_free_flags n k m Forbid}         & $\sigma$-typed $H$-free flags + completeness \\
\btt{generate_forbid_free_mul_theorems patN hostN k m Forbid} & $H$-free product expansions \\
\btt{expand_one_hfree_at n Forbid}                    & expand the unit over $H$-free flags \\
\btt{flag_expand_hfree n Forbid}                      & expand one flag over $H$-free flags \\
\bottomrule
\end{tabular}
\end{center}

\paragraph{File map.}
\begin{itemize}
  \item \btt{LeanFlagAlgebras/Forbid/Basic.lean} --- $\feq$, the vanishing lemma,
        and the restricted expansions (Section~\ref{sec:vanishing}).
  \item \btt{LeanFlagAlgebras/Flags/ForbidFreeGenerator.lean} --- the two
        generation commands, the \btt{toUnderlying} bridge, and the pruned
        generation + \btt{genFlagsHfree_toFinset_eq} (Sections~\ref{sec:generation},
        \ref{sec:pruned}).
  \item \btt{LeanFlagAlgebras/Flags/Densities/MulThmGenerator.lean} ---
        \btt{generate_forbid_free_mul_theorems}.
  \item \btt{LeanFlagAlgebras/API/Basic.lean} --- \btt{expand_one_hfree_at}.
  \item \btt{LeanFlagAlgebras/API/FlagExpand.lean} --- \btt{flag_expand_hfree}.
  \item \btt{LeanFlagAlgebras/Flagmatic/MantelHfree.lean} --- the worked example.
\end{itemize}

\end{document}
